\documentclass{article}
\date{}
\usepackage[synctex=auto]{../../apxproof}
\usepackage{hyperref}

\newtheoremrep{theorem}{Theorem}
\newtheoremrep{lemma}[theorem]{Lemma}
% A user environment wrapping a repeated theorem
\newenvironment{mylemma}[1][]{\begin{lemmarep}[#1]}{\end{lemmarep}}

\title{SyncTeX support}
\author{Author}

\begin{document}

\maketitle

With the \texttt{synctex} option, appendix material and repeated
theorems are re-read from the source files, so that SyncTeX points to
the right lines. The check script queries SyncTeX for the position of
the marked words on the appendix pages.

\section{Repeated Theorems}

\begin{theoremrep}[Named]
  \label{th:named}%
  This is the statement of the named theorem, marked quokka.
  It has an equation:
  \begin{equation}
    \label{eq:sum}
    \sum_{k=1}^\infty \frac1{k^2}=\frac{\pi^2}6
  \end{equation}
\end{theoremrep}

\begin{proof}
  This is the proof of the named theorem, marked wombat, and it refers
  to~(\ref{eq:sum-apx}).
\end{proof}

\begin{mylemma}[Wrapped]
  This lemma is wrapped in a user environment, marked numbat.
\end{mylemma}

\begin{proof}[Proof of the wrapped lemma]
  [This proof starts with a bracket and is marked bilby.]
\end{proof}

\begin{center}\begin{theoremrep}
  Centered statement, marked dingo.
\end{theoremrep}\end{center}

\begin{proof}
  Proof of the centered statement, marked echidna.
\end{proof}

\begin{theoremrep} Single-line statement, marked kookaburra. \end{theoremrep} Text after the end, marked platypus.

\begin{proof}
  Single-line proof, marked cassowary.
\end{proof}

\section{Material Sent to the Appendix}

\begin{toappendix}
  Appendix material, marked wallaby, with a proof:
  \begin{proof}[Nested]
    Nested proof, marked koala.
  \end{proof}
\end{toappendix}

Some main text.

\begin{toappendix}

\section{Own section, marked galah}
This appendix section replaces the automatic heading.
\end{toappendix}

Two cases found in real documents. First, an appendix block that consists of
a heading only: the paragraph that follows it comes from another block, and
must be treated as following a heading (no indentation, no page break) just
as when appendix material is copied.

\begin{toappendix}
  \subsection{Heading alone in its block}
\end{toappendix}
\begin{toappendix}
  Paragraph after the heading, marked kiwi.
\end{toappendix}

Second, an alignment ended in the block: the \verb|\end| of an environment must
remain expandable when \TeX{} scans the start of an alignment row, or an
empty row is added.

\begin{toappendix}
  \begin{center}
    \begin{tabular}{lr}
      \hline
      Marked & tenrec \\
      Second & row \\
      \hline
    \end{tabular}
  \end{center}
\end{toappendix}

Third, a repeated theorem stated inside an appendix block, which is then
just a theorem of the appendix.

\begin{toappendix}
  Before the nested theorem, marked lorikeet, on the same line.
  \begin{lemmarep}
    Nested lemma, marked pademelon.
  \end{lemmarep}
  \begin{proof}
    Its proof, marked possum.
  \end{proof}
  After the nested theorem, marked tuatara.
\end{toappendix}

\section{Included File}

\begin{lemmarep}
  Statement in the included file, marked emu.
\end{lemmarep}

\begin{proof}
  Proof in the included file, marked cockatoo.
\end{proof}


\end{document}
